\documentclass[../../main.tex]{subfiles}

\begin{document}
\begin{CJK*}{UTF8}{gbsn}

\section*{Exercise 3}

Suppose $(X_1,Y_1), \ldots, (X_{n+2},Y_{n+2}) \in \mathcal{X} \times \mathcal{Y}$ 
are exchangeable.
After observing $X_{n+1}$ and $X_{n+2}$, select $I \in \{1,2\}$, where
$I$ may depend on both $(X_{n+1}, X_{n+2})$. 
Given score $s : \mathcal{X} \times \mathcal{Y} \to \mathbb{R}$,
define a (full) conformal prediction set $\mathcal{C}(X_{n+I})$ such that 
\begin{equation*}
    P\{Y_{n+I} \in \mathcal{C}(X_{n+I})\} \geqslant 1-\alpha.
\end{equation*}

\paragraph{Solution}

The solution is to use a Benferroni correction for the two test points.

Specifically,
put $S_i = s(X_i,Y_i)$
and let $S_{(1)} \leqslant \cdots \leqslant S_{(n)}$
be the order statistics for $S_1, \ldots, S_n$.
Define

\begin{equation*}
    k=\left\lceil(n+1)\left(1-\frac{\alpha}{2}\right)\right\rceil,
\end{equation*}
\begin{equation*}
    \hat{q}=
    \begin{cases}
        S_{(k)}, & k \leqslant n,\\
        \infty, & k=n+1,
    \end{cases}
\end{equation*}
and 
\begin{equation*}
    \mathcal{C}(X_{n+I})
    =\curly{ y \in \mathcal{Y} \mid s(X_{n+I},y)\leqslant \hat{q}}.
\end{equation*}
Standard results (for example, for full conformal prediction,
Theorem 3.2 in the textbook) show that,
for $i \in \{1,2\}$, we have 

\begin{equation*}
    P\curly{Y_{n+i} \notin \mathcal{C}(X_{n+i}) } \leqslant \frac{\alpha}{2}.
\end{equation*}
Thus:

\begin{equation*}
    P \curly{Y_{n + I} \notin \mathcal{C}(X_{n+I})}
    = P \bracket{ \bigcup_{i=1}^2 \curly{Y_{n+i} \notin \mathcal{C}(X_{n+i})}
    } \leqslant 2 \times \frac{\alpha}{2} = \alpha.
\end{equation*}
This completes the proof.
////

\end{CJK*}
\end{document}
